% gradient-gapc.tex --- Section 5: Entropy Gradient Theorem and Gap C Resolution
% ASSERT_CONVENTION: natural_units=kB_equals_1, entropy_base=nats, state_normalization=Tr(rho)=1, sequential_product=a&b=sqrt(a)b*sqrt(a), von_neumann_entropy=S(rho)=-Tr(rho*ln*rho), information=I(B;M)=S(B)+S(M)-S(BM), free_energy=F=E-TS, clifford=Euclidean_Cl6_Lorentzian_Cl(d-1,1)

%----------------------------------------------------------------------
\section{Entropy Gradient Theorem and Gap~C
  Resolution}\label{sec:gradient-gapc}
\label{sec:gradient}% alias for cross-refs from Sec.~2
%----------------------------------------------------------------------

We now bring together the entropy dynamics of
Sec.~\ref{sec:entropy}, the Landauer bound of
Sec.~\ref{sec:landauer}, and the chirality--time-orientation link of
Sec.~\ref{sec:chirality} to establish the central results of this
paper.  The argument has four parts.  First, we prove that every
self-modeler requires an entropy gradient $S(t) < \Smax$---what we
call the \emph{entropy gradient theorem}---via three independent
routes (Sec.~\ref{sec:egt}).  Second, we identify two independent
chains that converge on time-orientation as a shared geometric
prerequisite (Sec.~\ref{sec:selection-chain}).  Third, we state the
narrowed Gap~C result as a formal theorem with epistemic labels
(Sec.~\ref{sec:gapc-theorem}).  Fourth, we present the updated master
theorem and nine-link chain with Link~6 annotated
(Secs.~\ref{sec:master-theorem}--\ref{sec:updated-chain}).

%----------------------------------------------------------------------
\subsection{Entropy gradient theorem}\label{sec:egt}
%----------------------------------------------------------------------

\begin{theorem}[Entropy gradient]\label{thm:entropy-gradient}
Let $(B,M,\mathbin{\&})$ be a self-modeling composite on a finite
SWAP lattice in thermal contact with an environment at
temperature~$T$, with $\dim\mathcal{H}_B = d_B$,
$\dim\mathcal{H}_M = d_M$, and faithful tracking
$\MI{B}{M} > 0$.  Then $S(t) < \Smax \equiv \ln(d_B\,d_M)$: the
composite must be out of thermal equilibrium, with entropy strictly
below its maximum.
\end{theorem}

\noindent
We prove this via three independent routes, each using a different
subset of assumptions.  Their convergence on the same conclusion
strengthens the result.

%......................................................................
\paragraph{Route~A (Direct chain).}
This route assembles the chain theorem of
Sec.~\ref{sec:chain} into a single implication.
\begin{enumerate}
  \item \textbf{Self-modeling $\boldsymbol{\Rightarrow}$ work.}
    Theorem~\ref{thm:landauer} gives
    $\Wmin = \kB T\,\MI{B}{M} > 0$ per update cycle.
    \hfill[\textsc{theorem}, Sec.~\ref{sec:landauer-bound}]

  \item \textbf{Work $\boldsymbol{\Rightarrow}$ non-equilibrium.}
    Performing work $W > 0$ requires free energy above the
    equilibrium value: $F(\rho_{BM}) > F_{\mathrm{eq}}$.
    Equivalently, the extractable work
    $\kB T\,D(\rho_{BM}\|\rho_{\mathrm{eq}}) > 0$ implies
    $\rho_{BM} \neq \rho_{\mathrm{eq}}$.
    \hfill[\textsc{standard physics}]

  \item \textbf{Non-equilibrium $\boldsymbol{\Rightarrow}$ entropy
    gradient.}
    A finite system out of equilibrium will equilibrate on a
    timescale set by the SWAP dynamics
    (Eq.~\eqref{eq:convergence}).  To sustain self-modeling beyond
    the exhaustion bound~\eqref{eq:exhaustion}, the system must
    access a low-entropy reservoir: $S(t) < \Smax$.
    \hfill[\textsc{physical argument} $+$ A3]
\end{enumerate}
The chain is:
\begin{equation}
  \MI{B}{M} > 0
  \;\xrightarrow{\;\text{Thm.~\ref{thm:landauer}}\;}
  W > 0
  \;\xrightarrow{\;\text{2nd Law}\;}
  F > F_{\mathrm{eq}}
  \;\xrightarrow{\;\text{A3}\;}
  S(t) < \Smax \,.
  \label{eq:route-a}
\end{equation}
Each link's output is the next link's input, and no link uses the
conclusion of a later link.

\paragraph{Route~B (Boundary argument).}
This route avoids the Landauer bound entirely, using only finitude and
monotonicity.

\begin{enumerate}
  \item \textbf{Entropy is bounded.}
    Finite-dimensional quantum mechanics gives
    $0 \leq S(\rho_{BM}) \leq \ln(d_B\,d_M) = \Smax$.
    \hfill[A1]

  \item \textbf{Entropy is non-decreasing.}
    Under repeated interaction with fresh maximally mixed bath qubits,
    the monotonicity chain~\eqref{eq:monotonicity} holds.  On a
    macroscopic lattice ($N \gg 1$), effective monotonicity holds for
    $t \ll t_{\mathrm{rec}} \sim e^{c\cdot 2^N}$
    (Proposition~\ref{prop:three-regimes}).
    \hfill[\textsc{theorem}, Sec.~\ref{sec:entropy-step}]

  \item \textbf{Self-modeling requires $S < \Smax$.}
    At $S = \Smax$ the state is maximally mixed, so $\MI{B}{M} = 0$
    (Eq.~\eqref{eq:eq-mutual}) and $\rhoexp = 0$.  Therefore
    $\rhoexp > 0$ implies $S(t) < \Smax$.
    \hfill[\textsc{definition}, Eq.~\eqref{eq:rho-def}]
\end{enumerate}
Combining: $S$ is bounded, non-decreasing, and currently below its
maximum.  Therefore for all $t_0 < t$,
\begin{equation}
  S(t_0) \;\leq\; S(t) \;<\; \Smax \,.
  \label{eq:route-b}
\end{equation}
The block possesses a low-entropy past.

\paragraph{Route~C ($\rhoexp$-selection).}
This route requires the fewest assumptions.  The contrapositive of
three established results gives:
\begin{equation}
  \rhoexp > 0
  \;\Longrightarrow\;
  \MI{B}{M} > 0
  \;\Longrightarrow\;
  \rho_{BM} \neq \rho_{\mathrm{eq}}
  \;\Longrightarrow\;
  S(t) < \Smax \,.
  \label{eq:route-c}
\end{equation}
Step~1 is the definition of $\rhoexp$
(Eq.~\eqref{eq:rho-def}: $\rhoexp = 0$ when $\MI{B}{M} = 0$).
Step~2 is the equilibrium result $\MI{B}{M}^{\mathrm{eq}} = 0$
(Eq.~\eqref{eq:eq-mutual}).  Step~3 is
$S(\rho_{\mathrm{eq}}) = \Smax$.  The $\rhoexp$-weighted ensemble is
therefore concentrated on blocks with $S < \Smax$: the experiential
measure sees only non-equilibrium blocks.

\medskip
\noindent
\textit{Comparison of routes.}---Table~\ref{tab:routes} compares the
three routes.  Route~A provides the strongest quantitative content
(the Landauer cost per cycle).  Route~B provides the simplest logic
(bounded monotone sequences).  Route~C requires the fewest
assumptions (only the definition of~$\rhoexp$ and the equilibrium
result).  Their convergence from different starting points makes the
entropy gradient theorem robust against the failure of any single
assumption.

\begin{table}[t]
  \centering
  \caption{Three routes to the entropy gradient theorem.  Each route
  proves $S(t) < \Smax$ using a different subset of assumptions.
  ``A$n$'' labels refer to the assumption register in
  Table~\ref{tab:assumptions}.}
  \label{tab:routes}
  \begin{tabular}{llll}
    \toprule
    & Route~A & Route~B & Route~C \\
    \midrule
    Key input
      & Landauer bound & Finitude $+$ monotonicity
      & $\rhoexp > 0 \Leftrightarrow I > 0$ \\
    Assumptions
      & A1, A2, A3, A4 & A1, A3 & A1, A5 \\
    Epistemic chain
      & Thm.\ $+$ std.\ phys.\ $+$ A3
      & Thm.\ (Lindblad) $+$ arithmetic
      & Def.\ $+$ thm.\ $+$ contrapositive \\
    Gives $W_{\min}$?
      & Yes & No & No \\
    Uses Landauer?
      & Yes & No & No \\
    \bottomrule
  \end{tabular}
\end{table}

%----------------------------------------------------------------------
\subsection{Two chains and their convergence}\label{sec:selection-chain}
%----------------------------------------------------------------------

The entropy gradient theorem and the chirality--time-orientation link
of Sec.~\ref{sec:chirality} provide two \emph{independent} chains of
reasoning.  We present each, then show that they converge on a common
geometric prerequisite: time-orientation.

%......................................................................
\paragraph{Chain~1 (all self-modelers).}
The entropy gradient theorem (Theorem~\ref{thm:entropy-gradient})
establishes that \emph{every} self-modeling composite requires an
entropy gradient, regardless of chirality or complexification.  The
logic, proved via three independent routes in Sec.~\ref{sec:egt}, is:
\begin{equation}
  \MI{B}{M} > 0
  \;\xrightarrow{\;\text{Landauer}\;}
  W \geq \kB T\,\MI{B}{M}
  \;\xrightarrow{\;\text{2nd Law}\;}
  \text{non-equilibrium}
  \;\xrightarrow{\phantom{\;\text{2nd Law}\;}}
  S(t) < \Smax \,.
  \label{eq:chain-1}
\end{equation}
This chain holds for all self-modelers on any block of structure
space.  It makes no reference to chirality, complexification, or
the Standard Model.  It uses only assumptions~A1--A3
(Table~\ref{tab:assumptions}).

The entropy gradient defines a temporal direction: the ``future'' is
the direction of entropy increase, and the ``past'' is the direction
of lower entropy.  On an emergent Lorentzian manifold (A6--A7), this
thermodynamic arrow requires a time-orientation---a continuous choice
of future-pointing timelike vector field---to distinguish the two
null half-cones at every point.  This last step uses A6--A7 and is
a physical identification (thermodynamic arrow $\to$ geometric
time-orientation), not a mathematical theorem; it is the standard
assumption that the emergent spacetime's causal structure is
compatible with its thermodynamics.

%......................................................................
\paragraph{Chain~2 (SM-like observers).}
Standard Model physics is defined by chiral fermions.  The chain
from complexification to chirality to time-orientation is:
\begin{equation}
  \text{SM-like observer}
  \;\xrightarrow{\;\text{def.}\;}
  \text{chirality}
  \;\xrightarrow{\;\text{\cite{Paper7}}\;}
  \text{complexification}
  \;\xrightarrow{\;\text{CHIR-01}\;}
  \text{time-orientation} \,.
  \label{eq:chain-2}
\end{equation}
Each arrow is a necessary condition, reading right-to-left:
\begin{itemize}
  \item \textbf{Chirality requires complexification.}
    On real $\Vhalf = \mathbb{O}^2$, the $\Cl{6}$ volume form
    $\omegasix$ has no eigenstates giving the left--right
    decomposition: the eigenvalue equation
    $i\,\omegasix\,\ket{\psi} = \pm\ket{\psi}$ requires complex
    coefficients.  Without complexification, chirality is not
    defined.
    \hfill[\textsc{algebra};~\cite{Paper7}]

  \item \textbf{Chirality requires time-orientation.}
    The three-consequence theorem
    (Theorem~\ref{thm:three-consequence}, consequence~(c))
    establishes that the physical Weyl decomposition
    $S = S_L \oplus S_R$ requires both internal chirality and
    a time-orientation for a globally consistent sign of the
    spacetime chirality operator $\Gammastar$
    (Proposition~\ref{prop:chirality-flip},
    Corollary~\ref{cor:projector-exchange}).
    \hfill[\textsc{geometry}; Sec.~\ref{sec:chirality}]
\end{itemize}
This chain uses assumptions~A6--A7 in addition to the algebraic
results of~\cite{Paper7}.

\paragraph{Logical status of the two chains.}
The direction of implication is critical.  The three-consequence
theorem proves \emph{chirality $\Rightarrow$ time-orientation is
required}.  Its valid contrapositive is \emph{no time-orientation
$\Rightarrow$ no chirality}.  Crucially, the \emph{converse}---no
chirality $\Rightarrow$ no time-orientation---does \emph{not}
follow.  A non-complexified block can lack chirality while still
possessing time-orientation from other sources.

This is why the two chains are independent.  Chain~1 establishes
entropy gradients (and hence a temporal direction) for all
self-modelers without invoking chirality.  Chain~2 shows that
SM-like observers additionally require chirality, which in turn
requires both complexification and time-orientation.  Neither chain
implies the other.

%......................................................................
\paragraph{Structural convergence.}
The two chains converge on time-orientation as a shared geometric
prerequisite:
\begin{itemize}
  \item \textbf{Chain~1:} the entropy gradient defines a temporal
    direction (low-entropy past $\to$ high-entropy future),
    which on a Lorentzian manifold corresponds to a
    time-orientation.
  \item \textbf{Chain~2:} chirality requires time-orientation for
    the Weyl decomposition
    (Theorem~\ref{thm:three-consequence}(c)).
\end{itemize}
The single algebraic choice $u \in S^6$ does double duty: it
determines the gauge group and chirality
(Paper~7~\cite{Paper7}) \emph{and} creates a requirement for the
same time-orientation that the thermodynamic arrow provides.  This
convergence is a structural consistency check: the framework's
particle physics content and its thermodynamic content both require
a geometric property (time-orientation) that the emergent spacetime
automatically provides.

It is \emph{not}, however, a proof that non-complexified blocks have
$\rhoexp = 0$.  A block without chirality can still possess
time-orientation (from Chain~1 alone), sustain an entropy gradient,
and support non-SM self-modelers.  The convergence constrains
SM-like observers; it does not exclude all observers from
non-complexified blocks.

Schematically, the two-chain structure is:
\begin{equation}
\boxed{
\begin{aligned}
  &\textbf{Chain 1:}\;
  \text{self-modeling}
  \;\Rightarrow\;
  S < \Smax
  \;\Rightarrow\;
  \text{time-orientation}
  \quad\text{(all self-modelers)}\\[4pt]
  &\textbf{Chain 2:}\;
  \text{SM-like}
  \;\Rightarrow\;
  \text{chirality}
  \;\Rightarrow\;
  \text{complexification} + \text{time-orientation}
  \quad\text{(SM observers)}
\end{aligned}
}
\label{eq:selection-chain}
\end{equation}

%----------------------------------------------------------------------
\subsection{Gap~C: narrowed}\label{sec:gapc-theorem}
%----------------------------------------------------------------------

We are now in a position to state the central result regarding
Gap~C\@.  The result is a \emph{narrowing}, not a full resolution:
complexification is necessary for SM-like observers, but we cannot
exclude non-SM self-modelers in non-complexified blocks.

\begin{theorem}[SM-observer complexification requirement]
\label{thm:complexification}
Let $\mathcal{S}$ be structure space with experiential measure
$\mu_{\rhoexp}$.  Let $\mathcal{B}$ be a block of $\mathcal{S}$
supporting a $C^*$-observer ($M_n(\mathbb{C})^{\mathrm{sa}}$)
probing~$\hthree$.  Suppose assumptions A1--A7 hold
\textup{(Table~\ref{tab:assumptions})}.  If
$\Vhalf(\mathcal{B})$ is not complexified
\textup{(}remains real $\mathbb{O}^2$\textup{)}, then:

\smallskip\noindent
\textup{(a)}~$\mathcal{B}$ has \textbf{no chiral spinors}: the
$\Cl{6}$ volume form $\omegasix$ has no eigenstates in real
$\Vhalf$ that give the $L/R$
decomposition.
\hfill\textup{[\textsc{algebra};~\cite{Paper7}]}

\smallskip\noindent
\textup{(b)}~$\mathcal{B}$ has \textbf{no SM gauge group}: without
complexification, the breaking chain
$\Spin{10} \to \PS \to \SMgauge$ is not
realized.
\hfill\textup{[\textsc{algebra};~\cite{Paper7}]}

\smallskip\noindent
\textup{(c)}~Chirality \textbf{requires time-orientation}: the
physical Weyl decomposition $S = S_L \oplus S_R$ requires a
time-orientation for a globally consistent sign of
$\Gammastar$.
\hfill\textup{[\textsc{geometry}; Sec.~\ref{sec:chirality}]}

\smallskip\noindent
\textup{(d)}~Self-modelers \textbf{independently require entropy
gradients}: $\MI{B}{M} > 0$ implies $S(t) < \Smax$ via
Theorem~\ref{thm:entropy-gradient}, with no reference to
chirality.
\hfill\textup{[\textsc{thermodynamics}; Sec.~\ref{sec:egt}]}

\smallskip\noindent
The convergence of \textup{(c)} and \textup{(d)} on
time-orientation means that within~$\hthree$, the particle physics
content \textup{(}gauge group, chirality\textup{)} and the
thermodynamic content \textup{(}entropy gradient, arrow of
time\textup{)} share a common geometric prerequisite.

\smallskip\noindent
Consequently, every block with SM-like observers \textup{(}chiral
fermions, SM gauge group\textup{)} has complexified
$\Vhalf^{\mathbb{C}} = S_{10}^+$.
\end{theorem}

\begin{proof}
Part~(a) is pure representation theory: on the real space
$\Vhalf = \mathbb{O}^2$, the operator $i\,\omegasix$ is not
defined ($i$ requires complex scalars), so the eigenvalue equation
$i\,\omegasix\ket{\psi} = \pm\ket{\psi}$ has no solutions.  The
left--right decomposition~\eqref{eq:chiral-decomp} therefore does
not exist.  This is proved in Ref.~\cite{Paper7}.

Part~(b) follows from~(a): the gauge symmetry breaking chain
$\Spin{10} \to \PS \to \SMgauge$ requires the complexified
representation $S_{10}^+$ for its chiral content.  Without
complexification, the $\Spin{10}$ structure acting on real
$\Vhalf$ does not decompose into SM representations.  This is
established in Ref.~\cite{Paper7}.

Part~(c) follows from
Theorem~\ref{thm:three-consequence}(c): the physical realization
of chirality as spacetime Weyl spinors requires both internal
chirality [from~(a)] and time-orientation
(Proposition~\ref{prop:chirality-flip}).  The direction of
implication is: chirality $\Rightarrow$ time-orientation required.
The valid contrapositive is: no time-orientation $\Rightarrow$ no
chirality.

Part~(d) is the entropy gradient theorem
(Theorem~\ref{thm:entropy-gradient}), proved via three independent
routes in Sec.~\ref{sec:egt}.  This part is logically independent
of~(a)--(c): it applies to all self-modelers regardless of
chirality or complexification.

The convergence statement follows from~(c) and~(d): both the
chirality requirement and the self-modeling requirement point to
time-orientation, but through independent logical paths.

The final statement (SM-like $\Rightarrow$ complexified) follows
from~(a) by contrapositive: no complexification $\Rightarrow$ no
chirality $\Rightarrow$ not SM-like.
\end{proof}

\begin{remark}[Selection, not algebraic forcing]
\label{rem:selection-not-forcing}
The complexification of $\Vhalf$ for SM-like observers is a
\emph{selection effect}, not an algebraic derivation.  Phase~22 of
this program established that four independent algebraic routes to
complexification all fail: the
$V_1 = \mathbb{R}\cdot E_{11}$ bottleneck in the Peirce
decomposition is a genuine obstruction~\cite{BGW2020}.
Theorem~\ref{thm:complexification} does not override this negative
result.  Non-complexified blocks exist in structure space; they are
not algebraically forbidden.  SM-like observers require
complexification, but the theorem does not establish
$\rhoexp = 0$ for all non-complexified blocks: a non-complexified
block could in principle support non-SM self-modelers with
$\rhoexp > 0$.
\end{remark}

\paragraph{What is proved and what is not.}
It is important to state the epistemic status precisely:
\begin{itemize}
  \item \textbf{Proved:} Complexification is necessary for SM-like
    observers (chiral fermions, SM gauge group).
  \item \textbf{Proved:} All self-modelers require entropy
    gradients, independent of chirality.
  \item \textbf{Proved:} Chirality and thermodynamics share
    time-orientation as a common geometric prerequisite.
  \item \textbf{Not proved:} $\rhoexp = 0$ for all
    non-complexified blocks.  A non-complexified block could have
    time-orientation, an entropy gradient, and non-SM self-modelers
    with $\rhoexp > 0$.
\end{itemize}
The last point does not undermine the program: the program explains
\emph{our} physics (SM particle content, chiral fermions,
Pati--Salam unification), and for that physics, complexification is
necessary.  Non-SM blocks, if they exist, are inaccessible under
L4---different structure-space points do not interact.

%......................................................................
\paragraph{Gap~C: narrowed.}
Paper~7's nine-link chain~\cite{Paper7} identifies Gap~C as the
complexification step $\Vhalf \to \Vhalf^{\mathbb{C}} = S_{10}^+$
(Link~L6).  Theorem~\ref{thm:complexification} \emph{narrows} this
gap: among all blocks supporting SM-like observers, $\Vhalf$ is
complexified.  The gap is not fully resolved because we cannot
exclude non-SM self-modelers in non-complexified blocks.  The
narrowing is via the experiential measure restricted to the SM
sector, not via algebraic necessity.  This is the appropriate level
of progress given Phase~22's proof that algebraic closure is
impossible~\cite{BGW2020}.

\paragraph{Status of the Past Hypothesis.}
The entropy gradient theorem
(Theorem~\ref{thm:entropy-gradient}) shows that self-modeling
\emph{requires} $S(t) < \Smax$.  Combined with the entropy
monotonicity of Sec.~\ref{sec:entropy}, this implies a low-entropy
boundary condition at earlier times
(Eq.~\eqref{eq:route-b}).  The Past Hypothesis---the
cosmological observation that the universe began in a state of
extraordinarily low entropy~\cite{Penrose1979,Albert2000}---is
thereby \emph{elevated} from ``observed feature of the universe'' to
``necessary condition for self-modeling.''  It is not
\emph{derived}: the entropy gradient theorem is silent on \emph{why}
the initial entropy was low.  It establishes only that it
\emph{must} have been low for blocks with $\rhoexp > 0$ to exist.

%----------------------------------------------------------------------
\subsection{The v7.0 master theorem}\label{sec:master-theorem}
%----------------------------------------------------------------------

We collect the results of this paper and its predecessors into a
single statement.

\begin{theorem}[Arrow of time and SM-observer complexification]
\label{thm:master}
Under assumptions A1--A7 \textup{(Table~\ref{tab:assumptions})},
the following hold for self-modeling composites
$(B,M,\mathbin{\&})$ on a finite SWAP lattice probing~$\hthree$:

\smallskip\noindent
\textup{(I)~\textbf{Entropy gradient.}}\enspace
Self-modelers require an entropy gradient: $S(t) < \Smax$.
This is established via three independent routes
\textup{(Table~\ref{tab:routes})}, each using a different assumption
subset.  This holds for all self-modelers regardless of chirality
or complexification.

\smallskip\noindent
\textup{(II)~\textbf{SM-observer complexification.}}\enspace
SM-like observers within~$\hthree$ require complexified $\Vhalf$.
Non-complexified blocks have no chiral spinors and no SM gauge
group \textup{(}Theorem~\ref{thm:complexification}(a)--(b)\textup{)}.

\smallskip\noindent
\textup{(III)~\textbf{Gap~C status: narrowed.}}\enspace
Complexification is necessary for SM-like observers
\textup{(}selection effect, not algebraic necessity\textup{)}.
The gap remains open for non-SM self-modelers: a non-complexified
block could in principle support observers with $\rhoexp > 0$ but
without chiral fermions or the SM gauge group.
\end{theorem}

\begin{corollary}[Chirality--thermodynamics nexus]
\label{cor:nexus}
The chirality of the Standard Model and the thermodynamic arrow of
time share a common geometric prerequisite:
\textbf{time-orientation} of the emergent spacetime.
Chirality requires time-orientation for the physical Weyl
decomposition
\textup{(Theorem~\ref{thm:three-consequence}(c))}.
The arrow of time requires time-orientation to define the direction
of entropy increase
\textup{(Theorem~\ref{thm:entropy-gradient})}.
Both are consequences of the self-modeling framework: the choice
$u \in S^6$ that determines particle physics
\textup{(}gauge group and chirality~\cite{Paper7}\textup{)}
simultaneously constrains spacetime geometry
\textup{(}time-orientation\textup{)}, and this time-orientation is
the same structure that supports the entropy gradient sustaining
self-modeling.  This convergence is structural, not causal: neither
chain implies the other.
\end{corollary}

%----------------------------------------------------------------------
\subsection{Updated nine-link chain}\label{sec:updated-chain}
%----------------------------------------------------------------------

Paper~7's nine-link chain from self-modeling to the chiral Standard
Model~\cite{Paper7} is updated with the results of this paper.
Table~\ref{tab:chain} shows each link, its status, and the
annotation from this work.

\begin{table*}[t]
  \centering
  \caption{Updated nine-link chain.  Links are numbered as in
  Paper~7~\cite{Paper7}.  ``Status'' indicates the epistemic
  classification; ``Annotation'' records changes from this paper and
  its predecessors~(Papers~5--7).  Link~L6 is narrowed via selection
  (Theorem~\ref{thm:complexification}); Link~L7 is updated with the
  three-consequence theorem
  (Theorem~\ref{thm:three-consequence}).}
  \label{tab:chain}
  \begin{tabular}{clll}
    \toprule
    Link & Statement & Status & Annotation \\
    \midrule
    L1 & Self-modeling composite $(B,M,\mathbin{\&})$ on $\hthree$
       & \textbf{Proved} \cite{Paper5}
       & --- \\
    L2 & Observer is $C^*$-algebra ($M_n(\mathbb{C})^{\mathrm{sa}}$)
       & \textbf{Proved} \cite{Paper5}
       & --- \\
    L3 & Faithful tracking requires local tomography
       & \textbf{Proved} \cite{Paper5}
       & --- \\
    L4 & $\hthree$ probed via rank-1 idempotent $e = E_{11}$
       & \textbf{Assumed} (Gap~B1)
       & --- \\
    L5 & Peirce decomposition: $V_1 \oplus \Vhalf \oplus V_0$
       & \textbf{Proved} \cite{Paper7}
       & --- \\
    L6 & $\Vhalf$ complexifies to $S_{10}^+$
       & \textbf{Narrowed via selection}
       & Thm.~\ref{thm:complexification}: necessary
         for SM-like; open for non-SM \\
    L7 & $u \in S^6$ determines gauge $+$ chirality
         $+$ time-orient.
       & \textbf{Proved $+$ updated}
       & Thm.~\ref{thm:three-consequence}: third
         consequence~(c) \\
    L8 & $\Spin{10} \to \PS \to \SMgauge$
       & \textbf{Proved} \cite{Paper7}
       & --- \\
    L9 & Chiral SM fermions from $\Cl{6}$ action
       & \textbf{Proved} \cite{Paper7}
       & --- \\
    \bottomrule
  \end{tabular}
\end{table*}

\noindent
Of the nine links, six are proved algebraically (L1--L3, L5, L8,
L9), one is assumed (L4, Gap~B1), one is updated with a new
geometric consequence (L7), and one is narrowed via selection (L6,
this paper): necessary for SM-like observers, open for non-SM
self-modelers.

%----------------------------------------------------------------------
\subsection{Assumption register}\label{sec:assumption-register}
%----------------------------------------------------------------------

Table~\ref{tab:assumptions} lists every assumption entering the
master theorem (Theorem~\ref{thm:master}), ordered by severity.
Assumption~A3 is the strongest: it connects local thermodynamics to
cosmological initial conditions and is the only assumption that, if
it fails, weakens the ``low-entropy past'' conclusion of
Route~A\@.  Assumptions~A6--A7 are needed only for the
SM-observer complexification requirement [part~(II) of the master
theorem]; the entropy gradient theorem [part~(I)] survives their
failure.

\begin{table*}[t]
  \centering
  \caption{Assumption register for the master theorem
  (Theorem~\ref{thm:master}).  Assumptions are ordered by decreasing
  severity (A3 strongest, A1/A5 weakest).  ``Used in'' refers to the
  parts of the master theorem.}
  \label{tab:assumptions}
  \begin{tabular}{cllll}
    \toprule
    ID & Assumption & Source & Status
       & Used in \\
    \midrule
    A1 & Finite-dimensional QM
       & Paper~5~\cite{Paper5}
       & Axiom
       & (I) all routes; (II) all steps \\
    A2 & Thermal contact at temperature $T$
       & Standard~\cite{ReebWolf2014}
       & Standard physics
       & (I) Route~A, Link~1 \\
    A3 & Finite system equilibrates on finite timescales
       & Stat.\ mech.
       & Physical argument
       & (I) Routes~A,B; Link~3 \\
    A4 & SWAP lattice dynamics $H = JF$
       & Paper~6~\cite{Paper6}
       & Model choice
       & (I) timescales only \\
    A5 & $\rhoexp = \MI{B}{M}(1 - \MI{B}{M}/\SvN{\rho_B})$
       & Paper~5~\cite{Paper5}
       & Axiom
       & (I) Route~C; (II) step~(d) \\
    A6 & Continuum limit yields smooth manifold
       & Paper~6~\cite{Paper6}
       & Assumption
       & (II) steps~(a)--(c) \\
    A7 & Lorentzian signature $(d{-}1,1)$
       & Paper~6~\cite{Paper6}
       & Physical
       & (II) steps~(a)--(c) \\
    \bottomrule
  \end{tabular}
\end{table*}

%----------------------------------------------------------------------
\subsection{Explicit non-claims}\label{sec:explicit-nonclaims}
%----------------------------------------------------------------------

To prevent overclaiming, we state what the results of this section
do \emph{not} establish.

\begin{enumerate}
  \item \textbf{Gap~C is not closed algebraically.}
    The algebraic obstruction ($V_1 = \mathbb{R}\cdot E_{11}$) found
    in Phase~22~\cite{BGW2020} remains genuine.
    Theorem~\ref{thm:complexification} narrows Gap~C via the
    experiential measure restricted to SM-like observers, not by
    overriding the algebraic result.

  \item \textbf{$\rhoexp = 0$ is not proved for all non-complexified
    blocks.}
    A non-complexified block could possess time-orientation, sustain
    an entropy gradient, and support non-SM self-modelers with
    $\rhoexp > 0$.  Complexification is necessary for SM-like
    observers, not for all observers.

  \item \textbf{No chirality does not imply no time-orientation.}
    The three-consequence theorem proves chirality $\Rightarrow$
    time-orientation.  The valid contrapositive is
    no time-orientation $\Rightarrow$ no chirality.  The
    converse---no chirality $\Rightarrow$ no time-orientation---is
    logically invalid and is not used in any argument of this paper.

  \item \textbf{This does not derive the Past Hypothesis.}
    The entropy gradient theorem shows that $S(t) < \Smax$ is
    \emph{necessary} for self-modeling.  It does not explain
    \emph{why} the universe started with low entropy.

  \item \textbf{The Past Hypothesis is elevated, not derived.}
    Its status changes from ``observed cosmological fact'' to
    ``necessary condition for self-modeling''---an
    observer-selection argument, not a first-principles
    derivation~\cite{Penrose1979,Albert2000}.

  \item \textbf{``Narrowed via selection'' is weaker than ``resolved
    via selection,'' which is weaker than ``proved algebraically.''}
    An algebraic proof of complexification would close Gap~C
    unconditionally.  A full selection resolution would show
    $\rhoexp = 0$ for all non-complexified blocks.  The narrowing
    establishes necessity for SM-like observers only.
\end{enumerate}

\medskip
With Gap~C narrowed and the master theorem stated,
we turn in Sec.~\ref{sec:predictions} to the testable consequences
of the framework.
